\section{Motivation}

The field of modern robotics has achieved unprecedented success in dexterous manipulation, largely driven by breakthroughs in imitation learning, particularly the development of Vision-Language-Action (VLA) models. These architectures have demonstrated remarkable capabilities in semantic reasoning, zero-shot generalization, and the execution of complex tasks in unstructured environments. However, the performance of these models is fundamentally constrained by the availability of high-quality human demonstration data. To capture smooth, human-like motions, researchers have to teleoperate physical systems using hardware interfaces that often bottleneck the data collection process.

Teleoperation of robot manipulators traditionally relies on teach pendants, joysticks, or spacemouse devices that map low-dimensional human input to the robot's joint or task space. These interfaces are unintuitive for 6-DOF control: operators must mentally decompose desired end-effector motions into separate translational and rotational commands, often through menus or mode switches. This disconnect between the operator's spatial intent and the input device limits both the speed and precision of teleoperation.

Modern smartphones present a compelling alternative. Equipped with inertial measurement units (IMU), cameras, and in some models LiDAR sensors, devices such as the iPhone can estimate their own 6-DOF pose in real time via visual-inertial odometry (VIO). Since humans naturally manipulate phones in 3D space---translating, rotating, and tilting them with high dexterity---the phone itself becomes an intuitive spatial controller that directly maps to the robot's task space.